%%=============================================================================
%% Voorwoord
%%=============================================================================

\chapter*{\IfLanguageName{dutch}{Woord vooraf}{Preface}}%
\label{ch:voorwoord}

Deze bachelorproef werd geschreven in het kader van mijn opleiding Toegepaste Informatica (afstudeerrichting Full Stack Development) aan de Hogeschool Gent.

De keuze voor dit onderwerp kwam voort uit mijn frustratie in het dagelijkse leven over het aantal gebruiksonvriendelijke applicaties. Ik merkte vaak dat formulieren onduidelijk zijn, waardoor ik regelmatig mijn grootouders, ouders en vrienden moest helpen bij het invullen ervan. Dit zette mij aan het denken: hoe kan dit probleem worden aangepakt?

Tijdens mijn opleiding stelde ik vast dat gebruiksvriendelijkheid bij applicatieontwikkeling niet altijd de hoogste prioriteit krijgt, of dat er onvoldoende budget voorzien wordt voor usabilitytesten. Daarom wilde ik onderzoeken hoe deze testen het best en efficiënt kunnen worden uitgevoerd. Op die manier hoop ik niet alleen anderen te kunnen helpen, maar ook toekomstige frustraties bij mezelf te beperken.

Tijdens het schrijven van deze bachelorproef heb ik veel bijgeleerd over wanneer en hoe usabilitytesten het meest effectief zijn. Ik hoop dat deze inzichten in de toekomst zowel voor mezelf als voor anderen een meerwaarde kunnen betekenen.

Graag wil ik mijn promotor, dhr.\ P. Beke, bedanken voor de begeleiding en waardevolle feedback gedurende dit proces. Daarnaast wil ik ook mijn copromotor, Robbe Doeven, en Delaware bedanken voor hun ondersteuning en het vertrouwen in mijn onderzoek.

Verder wil ik alle vrijwilligers bedanken die hebben deelgenomen aan mijn onderzoek. Zonder hun bijdrage zou deze bachelorproef niet mogelijk zijn geweest.

Tot slot wil ik mijn ouders bedanken voor de kansen die zij mij hebben gegeven en voor hun voortdurende steun gedurende mijn studies. Ook wil ik mijn vrienden bedanken, en in het bijzonder mijn huisgenoten Tim en Julie, voor hun steun en luisterend oor.

\bigskip
\noindent Gent, 3 mei 2026

\medskip
\noindent Fien Teugels